% Nam slide tom tat Sang kien cong nghe VDT 2026.
% Bien dich: pdflatex slides.tex (chay hai lan neu can).
\documentclass[aspectratio=169,11pt]{beamer}
\usepackage[utf8]{inputenc}
\usepackage[T5]{fontenc}
\usepackage[vietnamese]{babel}
\usepackage{lmodern}
\usepackage{booktabs}
\usepackage{tikz}
\usetikzlibrary{arrows.meta,positioning}
\usetheme{Madrid}
\usecolortheme{default}
\definecolor{vdtblue}{RGB}{20,78,128}
\setbeamercolor{structure}{fg=vdtblue}
\setbeamertemplate{navigation symbols}{}
\setbeamertemplate{footline}[frame number]
\newcommand{\code}[1]{\texttt{#1}}

\begin{document}

\begin{frame}[plain]
\centering
\includegraphics[page=1,trim=175bp 640bp 175bp 45bp,clip,width=3.1cm]{vdt_report_sang_kien_cong_nghe.pdf}\\[0.1cm]
{\small\bfseries TẬP ĐOÀN CÔNG NGHIỆP - VIỄN THÔNG QUÂN ĐỘI}\\[0.5cm]
{\Large\bfseries BÁO CÁO SÁNG KIẾN CÔNG NGHỆ}\\[0.35cm]
{\LARGE\bfseries Mô phỏng và trực quan hóa}\\[0.08cm]
{\LARGE\bfseries chỉ số vận hành UPF 5G}\\[0.35cm]
{\normalsize\bfseries Họ và tên: Trần Phúc Long}\\[0.12cm]
{\small\bfseries Email: longtp10@viettel.com.vn}\\[0.12cm]
{\small\bfseries Chương trình Viettel Digital Talent 2026}\\[0.12cm]
{\small\bfseries Lĩnh vực: Software Engineer}
\vfill
\begin{tabular}{ll}
\textbf{Mentor:} & Nguyễn Ngọc Dũng \\
\textbf{Đơn vị:} & Viettel High Tech
\end{tabular}
\end{frame}

\begin{frame}{Giải pháp và kiến trúc}
\centering
\begin{tikzpicture}[node distance=0.65cm, box/.style={draw=vdtblue,very thick,rounded corners,align=center,minimum width=3.3cm,minimum height=1.45cm}, arr/.style={-{Latex},very thick,vdtblue}]
\node[box] (go) {Go mock UPF\\5 metric families\\\code{/metrics}, \code{/mode}};
\node[box,right=of go] (prom) {Prometheus\\scrape mỗi 5 giây\\lưu chuỗi thời gian};
\node[box,right=of prom] (graf) {Grafana\\dashboard 4 panel\\refresh mỗi 5 giây};
\draw[arr] (go) -- (prom);
\draw[arr] (prom) -- (graf);
\end{tikzpicture}
\vspace{0.7cm}
\begin{itemize}
\item Docker Compose dựng toàn bộ ba dịch vụ bằng \code{docker compose up --build}.
\item Dashboard và datasource được nạp tự động từ cấu hình trong repository.
\item API chuyển \code{normal} / \code{degraded} để trình diễn sự cố.
\end{itemize}
\end{frame}

\begin{frame}{Phần đã triển khai}
\begin{columns}[T]
\begin{column}{0.52\textwidth}
\begin{block}{Nguồn dữ liệu xác định}
Mỗi giây cập nhật theo nhịp \(w=t\bmod 10\). Không dùng random; cùng kịch bản tạo lại được cùng mẫu dữ liệu.
\end{block}
\begin{block}{Bất biến dữ liệu}
\(\Delta RX=\Delta TX+\Delta Drop\). Counter chỉ tăng; session và lcore là Gauge.
\end{block}
\end{column}
\begin{column}{0.45\textwidth}
\begin{block}{Dashboard}
\begin{enumerate}
\item Packet throughput RX/TX
\item Packet drop rate
\item Active sessions
\item Per-lcore utilization (0--3)
\end{enumerate}
\end{block}
\end{column}
\end{columns}
\vspace{0.25cm}
\centering\small Mã nguồn Go, kiểm thử, Compose, cấu hình Prometheus và Grafana đã có trong repository.
\end{frame}

\begin{frame}{Kết quả triển khai và kiểm chứng}
\small
\begin{table}
\centering
\begin{tabular}{lcc}
\toprule
\textbf{Chỉ số mô phỏng mỗi giây} & \textbf{Normal} & \textbf{Degraded} \\
\midrule
Gói nhận & 1.200--1.335 & 520--592 \\
Gói bị loại & 2--4 & 75--93 \\
Tải lcore 0 & 32--36,5\% & 78--82,5\% \\
\bottomrule
\end{tabular}
\end{table}
\vspace{0.2cm}
\begin{itemize}
\item \code{go test ./...}: đạt; kiểm tra bất biến Counter, mode switch, lcore và mode lỗi.
\item \code{docker compose config --quiet}: đạt; cấu hình triển khai hợp lệ.
\item Các dải số trên suy ra từ công thức simulator, \textbf{không phải đo UPF thật}.
\end{itemize}
\end{frame}

\begin{frame}{Ứng dụng thực tế và trình diễn}
\begin{block}{Giá trị hiện tại}
Demo giám sát UPF, đào tạo cách đọc metric, thử dashboard và quy trình triển khai trước khi nối nguồn dữ liệu thật.
\end{block}
\begin{enumerate}
\item Mở \code{localhost:3000/d/mock-upf-demo/mock-upf-overview}; quan sát bốn panel ở chế độ normal.
\item Gọi \code{POST /mode?mode=degraded} trên cổng 8080; đợi khoảng 20 giây để thấy throughput giảm, drop và tải lcore tăng.
\item Gọi \code{POST /mode?mode=normal} để khôi phục; đối chiếu \code{localhost:9090/targets}.
\end{enumerate}
\vspace{0.2cm}
\small\textbf{Bước tiếp theo:} kết nối metric UPF thật, thêm cảnh báo và đánh giá end-to-end bằng dữ liệu đo.
\end{frame}

\end{document}
